%% Referee report on "Multiplicity one for SL(2)", Draft 3 (D. Joyner)
\documentclass[12pt,letterpaper]{article}
\usepackage[american]{babel}
\usepackage[T1]{fontenc}
\usepackage{amsmath,amssymb,amsthm}
\usepackage{mathptmx}
\usepackage{courier}
\usepackage[letterpaper,margin=1in]{geometry}
\usepackage{microtype}
\usepackage{enumitem}
\usepackage{booktabs}
\usepackage{longtable}
\usepackage[hidelinks]{hyperref}

\newcommand{\tr}{\operatorname{tr}}
\newcommand{\Ad}{\operatorname{Ad}}
\DeclareMathOperator{\sgn}{sgn}
\newcommand{\A}{\mathbb{A}}
\newcommand{\R}{\mathbb{R}}
\newcommand{\Q}{\mathbb{Q}}
\newcommand{\C}{\mathbb{C}}
\newcommand{\Z}{\mathbb{Z}}
\newcommand{\ov}[1]{\overline{#1}}
\newcommand{\lref}[1]{\textup{(l.\,#1)}}

\theoremstyle{plain}
\newtheorem*{propR}{Proposition R}
\newtheorem*{corR}{Corollary R}

% Suggested replacement text is set off like this.
\newenvironment{fix}[1][Suggested text]%
  {\begin{quote}\small\noindent\textbf{#1.}\ }%
  {\end{quote}}

\setlist[enumerate]{itemsep=3pt,topsep=4pt}
\setlist[itemize]{itemsep=2pt,topsep=4pt}

\title{Referee report\\[1ex]
\large D.~Joyner, \emph{Multiplicity one for $SL(2)$}, Draft~3 (November 2026)\\
file \texttt{multiplicity-one-SL2\_draft3.tex}}
\author{Claude (Anthropic)}
\date{November 23, 2026}

\begin{document}
\maketitle

\section{Summary and recommendation}

The draft is honest about its status, and most of what it labels as proved
is proved. The proof of Theorem~7.9 (\texttt{thm:main}) is logically
correct as a deduction from Hypotheses LH and GH. The local torus
computations of \S5.3, the real elliptic matchings of Lemma~5.29
(\texttt{lem:realds}), Lemma~7.12 (\texttt{lem:A00}), and
Proposition~7.15 (\texttt{prop:LH-spherical}) all check.

There is, however, one finding that changes the status of the main result.

\begin{quote}
\textbf{The hypotheses of the main theorem can never hold.} For every
cuspidal $\pi$ of $SL(2,\A)$ and every finite $S\supseteq S(\pi)$, there
is \emph{no} genuine cuspidal representation $\ov\pi$ whose components
outside $S$ correspond to those of $\pi$ under the Satake correspondence of
Definition~5.21 (\texttt{def:satake-corr}). Conclusion (b) of the theorem
is therefore false for every $\pi$, so Hypotheses LH and GH are jointly
inconsistent and the theorem is vacuous. See M1 below, where the argument
is written out.
\end{quote}

The cause is the choice of unramified correspondence. The draft matches
unramified principal series ``pointwise'' on the torus (the \emph{odd}
correspondence), whereas the correspondence that exists globally, through
Waldspurger and restriction from $PGL(2)$, is the squaring (\emph{even})
one. The draft's own archimedean matchings are of the even type, so the
two ends of the comparison do not fit together.

Six further issues are substantive (M2--M7): Hypothesis GH depends on
unstated local choices; Hypothesis LH as stated fails for the real elliptic
classes; the supercuspidal relation of Remark~2026.1 is false, not merely
unverified, when $q\equiv3\pmod4$; the formula for the weighted orbital
integral at central elements diverges; the recalled trace formula omits
terms; and three statements in proved or explanatory text are wrong
(Remark~2026.11 and Lemmas 3.2 and 3.4).

\textbf{Recommendation: major revision.} The local results are worth
keeping. Theorem~7.9 should not be presented as a conditional proof of
multiplicity one in its present form. Section~\ref{sec:options} lists the
options I see.

\section{What was checked, and how}

\begin{itemize}
\item \emph{Read in full:} all 5{,}253 source lines. The file compiles
with \texttt{pdflatex} (72 pages, no errors, no undefined references or
citations).
\item \emph{Re-derived by hand:} the cocycle identities of \S2.1,
including Lemma~2.3(h); Lemmas 5.3, 5.4, 5.8, 5.15 and the two
$\Q_3$ examples; the Haar measure formula $dg=\delta_B(a)\,dk\,da\,dn$; the
character formulas (5.17) and the table in the proof of
Lemma~5.29; Remark~5.27(2); the Jacquet module computation in
Lemma~7.12(ii); every step of the proof of Theorem~7.9.
\item \emph{Checked by exact computation} (script
\texttt{referee\_checks.py}, supplied with this report; Hilbert symbols over
$\Q_p$, $p=3,5,7$, on random elements of $SL(2,\Z)$): that
$k\mapsto(k,s_v(k))$ is a homomorphism on $SL(2,\Z_p)$; Lemma~2.3(h); and
the counterexample in M4.
\item \emph{Checked against the web:} the statement and bibliographic data
of Chai--Zhang and of Zhang's addendum; the titles and status of W.-W.~Li's
papers on the metaplectic trace formula.
\item \emph{Not consulted directly:} Waldspurger \cite{Wa1,Wa2},
Gelbart--Jacquet, Jacquet--Shalika, Arthur, Gelbart's lecture notes, Adams,
Manderscheid, Savin, Shokranian. Where I use results from these (notably in
M1 and M5) I quote them from memory. They are standard, but the author
should confirm each one before relying on my argument. Bibliographic
corrections marked ``from memory'' in \S\ref{sec:bib} carry the same
caveat.
\end{itemize}

An AI referee report is not an independent certification. M1 in particular
should be read by a specialist.

\section{Major issues}

\subsection*{M1. The hypotheses of Theorem 7.9 are mutually inconsistent}

\paragraph{The local computation.} Let $v$ be odd, $F=k_v$, and let
$\psi$ be an additive character of $F$ of conductor ${\cal O}$. Put
$\chi_\psi(x)=\gamma(q)/\gamma(xq)$, the function $\omega$ of \S5.7.2, so
that $\chi_\psi(xy)=(x,y)\chi_\psi(x)\chi_\psi(y)$, $\chi_\psi=1$ on
${\cal O}^\times$ and $\chi_\psi(x^2)=1$. By Lemma~5.4 every unramified
genuine character of $\ov{A(F)}$ is
$(a(x),\zeta)\mapsto\zeta\,\chi_\psi(x)\mu(x)$ for a unique unramified
character $\mu$ of $F^\times$. Suppose this character is
$\ov\chi_\epsilon$ and corresponds to $\chi=\chi_\vartheta$ in the sense of
Definition~5.5. Writing $\ov\psi(x)=\ov\chi(a(x),1)$, the cocycle relation
gives $\ov\psi(\varpi)^2=(\varpi,\varpi)\,\chi_1(\varpi^2)$ and
$\chi_\psi(\varpi)^2=(\varpi,\varpi)$, hence
\begin{equation}\label{eq:square}
\mu(\varpi)^2=\chi_1(\varpi^2)=\chi(a(\varpi))^2,
\qquad\text{i.e.}\qquad \mu(\varpi)=\pm\,\chi(a(\varpi)).
\end{equation}
The sign depends on $\varpi$, $\eta_\varpi$ and $\psi$ (see M2). So the
draft's correspondence is $t'=\pm t$, where $t=\chi(a(\varpi))$ is the
Satake parameter on $SL(2)$ and $t'=\mu(\varpi)$ the one on the cover.

\paragraph{What the global theory gives.} For a genuine unramified
principal series $I(\chi_\psi\mu)$, Waldspurger's lift to $PGL(2,F)$ is
$\pi(\mu,\mu^{-1})$, whose restriction to $SL(2,F)$ has torus character
$a(x)\mapsto\mu(x)^2$. On dual groups this is the projection
$SL(2,\C)\to PGL(2,\C)$. The functorial relation is therefore
$t=t'^{\,2}$, not $t=\pm t'$.

\begin{propR}
Let $\pi$ be a cuspidal automorphic representation of $SL(2,\A)$ and
$S\supseteq S(\pi)$ a finite set of places. There is no genuine cuspidal
automorphic representation $\ov\pi'$ of $\ov{SL(2,\A)}$ such that
$\ov\pi'_v$ corresponds to $\pi_v$ under the Satake correspondence of
Definition~5.21 for every $v\notin S$. This holds for every choice of the
uniformizers $\varpi_v$ and of the constants $\eta_{\varpi_v}$.
\end{propR}

\begin{proof}
Suppose $\ov\pi'$ exists. By Lemma~7.12 it lies in ${\cal A}_{00}$. Fix a
nontrivial character $\psi$ of $k\backslash\A$. By Waldspurger
\cite{Wa1,Wa2}, $\sigma=\mathrm{Wald}_\psi(\ov\pi')$ is a cuspidal
automorphic representation of $PGL(2,\A)$, and for almost all $v$, if
$\ov\pi'_v\subset I(\chi_{\psi_v}\mu_v)$ then
$\sigma_v=\pi(\mu_v,\mu_v^{-1})$. By Labesse--Langlands, $\pi$ occurs in
the restriction of a cuspidal $\Pi$ of $GL(2,\A)$; write
$\Pi_v=I(\mu_{1,v},\mu_{2,v})$ and $t_v=(\mu_{1,v}/\mu_{2,v})(\varpi_v)$ at
unramified places. By \eqref{eq:square}, $\mu_v(\varpi_v)^2=t_v^2$ for
almost all $v$. The Satake parameters of the adjoint lifts are then
\[
\Ad(\Pi)_v:\ \{t_v,\,1,\,t_v^{-1}\},\qquad
\Ad(\sigma)_v:\ \{t_v^2,\,1,\,t_v^{-2}\}.
\]
Multiplying out, the nine parameters of $\Ad(\Pi)_v\otimes\Ad(\Pi)_v$ are
$\{t_v^{2},1,t_v^{-2}\}\cup\{t_v,1,t_v^{-1}\}\cup\{t_v,1,t_v^{-1}\}$.
Hence, for a sufficiently large finite set $S'$,
\[
L^{S'}\!\bigl(s,\Ad(\Pi)\times\Ad(\Pi)\bigr)
=L^{S'}\!\bigl(s,\Ad(\sigma)\bigr)\,L^{S'}\!\bigl(s,\Ad(\Pi)\bigr)^{2}.
\]
By Gelbart--Jacquet \cite{GJ78}, $\Ad(\Pi)$ is a self-dual isobaric
automorphic representation of $GL(3,\A)$ with unitary cuspidal
constituents, so the left side has a pole at $s=1$ (Jacquet--Shalika). On
the right, $L^{S'}(s,\Ad\sigma)=L^{S'}(s,\sigma\times\sigma^\vee)/\zeta_k^{S'}(s)$
is holomorphic at $s=1$ because $\sigma$ is cuspidal, and likewise
$L^{S'}(s,\Ad\Pi)$. This is a contradiction.
\end{proof}

\begin{corR}
No cuspidal $\pi$ satisfies the hypotheses of Theorem~7.9. If
Hypothesis~LH holds at the places of $S(\pi)$, then Hypothesis~GH is false
for $\nu=\nu_\infty(\pi)$.
\end{corR}

Indeed, the proof of Theorem~7.9 derives the existence of $\ov\pi'$ from
the hypotheses, and that derivation is correct.

\paragraph{Symptoms visible inside the draft.}
\begin{enumerate}
\item \emph{Reducibility and unitarity are not preserved.} With
$|\chi(a(x))|=|\ov\chi(a(x),1)|=|x|^s$, the genuine principal series is
unitary for $s<1/2$ and reducible at $s=1/2$ (Lemmas 3.7 and 3.8), while
$I_A(\chi)$ on $SL(2,F)$ is unitary for $s<1$ and reducible at $s=1$.
Lemma~7.12(ii) uses exactly this: it matches the even Weil representation
with an \emph{irreducible complementary series} of $SL(2,F)$ at $s=1/2$.
Under $t=t'^{\,2}$ it matches the trivial representation, as it should.
\item \emph{The archimedean matchings are of the other type.}
Definition~5.26 relates $\Phi_{\ov f}(\ov r(2\theta))$ to
$\Phi_f(r(\theta))$, that is, it uses the squaring map
$h\mapsto(h,1)^2$, and Lemma~5.29(a) sends weight $k/2$ to
$m=\pm(k-2)$, weight $k-1$. That is the Shimura correspondence. The lemma
on non-tempered representations
($\Theta_{\ov\pi_{+,nt}}+\Theta_{\ov\pi_{-,nt}}=\Theta_{\pi_{nt}}$) is
also of even type. The hyperbolic matchings (Definitions 5.7, 5.9, 5.16, 5.21)
are of odd type. Hypothesis GH combines the two.
\item \emph{The running example.} For $\Delta$ the draft finds
$\ov\pi_\infty=\ov\pi^{\pm}_{13/2}$. The genuine form that actually exists
is the weight $13/2$ form attached to $\Delta$ by Shimura--Kohnen. Its
unramified parameters satisfy $t_p=t_p'^{\,2}$, not $t_p=\pm t'_p$.
\end{enumerate}

\paragraph{Literature.} The modern theory supports the same conclusion.
W.-W.~Li has stabilized the full trace formula for $Mp(2n)$; the elliptic
endoscopic groups are products of split odd orthogonal groups, so for
$n=1$ the comparison is with $SO(3)=PGL(2)$ (and its inner form), not with
$SL(2)$. This is the same phenomenon the draft already notes for Hecke
algebras in Remark~5.31 (\texttt{rem:savin-caveat}). The literature review
should cite this work (see \S\ref{sec:bib}).

\subsection*{M2. The Satake correspondence depends on arbitrary local choices}

Independently of M1, Hypothesis GH is not well posed as written. The
bijection $\chi\mapsto\ov\chi$ of Definition~5.5 depends on $\varpi$ and
$\eta_\varpi$, as the draft says, but the dependence is not harmless:
\begin{itemize}
\item replacing $\eta_\varpi$ by $-\eta_\varpi$ (possible when
$q\equiv3\pmod4$) composes the bijection with the twist by the unramified
quadratic character $a(x)\mapsto(-1)^{\nu(x)}$;
\item replacing $\varpi$ by $u\varpi$ with $u$ a non-square unit does the
same, for every odd $q$. (With $\varpi'=u\varpi$ one finds
$\vartheta'=\chi_1(u)\vartheta$ and $\epsilon'=-\chi_1(u)\epsilon$.)
\end{itemize}
So at each odd place there are two ``Satake correspondences,'' differing by
an unramified quadratic twist, and Definition~7.2 uses one at every
$v\notin S$ without saying which. Hypothesis GH is a different statement
for each of the uncountably many families of choices. By Proposition~R it
is false for all of them, but a reader should not have to discover that:
the dependence should be stated where GH is stated. The even
correspondence $\chi(a(x))=\mu(x)^2$ is canonical: changing $\psi$ to
$\psi^a$ multiplies $\mu$ by $(a,\cdot)$ and leaves $\mu^2$ unchanged.

A related omission: the draft never checks that
$\chi\mapsto\ov\chi$ commutes with the Weyl group, although \S5.3.4 and
Proposition~5.12 use it. It does:

\begin{fix}[Suggested lemma, after Definition 5.5]
If $\ov\chi$ corresponds to $\chi$, then $\ov\chi^{\,w}$ corresponds to
$\chi^w=\chi^{-1}$. \emph{Proof.} By Lemma~5.3(b),
$\ov\chi^{\,w}(a(x),1)=\ov\chi(a(x^{-1}),1)$, which restricts to
$\chi_1^{-1}$ on $A_2(F)$ and has parameter
$\epsilon''=(\varpi,\varpi)\ov\chi(a(\varpi^{-1}),1)
=\epsilon\,\chi_1(\varpi^2)^{-1}=(\varpi,\varpi)\epsilon^{-1}
=(\varpi,\varpi)\eta_\varpi^{-1}\vartheta^{-1}$. Since
$\eta_\varpi^2=(\varpi,\varpi)$, this equals $\eta_\varpi\vartheta^{-1}$,
the parameter of the genuine character corresponding to $\chi^{-1}$.
\end{fix}

\subsection*{M3. Hypothesis LH fails for the real elliptic classes}

Section~7.4, item~5, says that (LH1)--(LH3) are ``not carried out here''
for the classes of \S5.5. In fact (LH1) and (LH2), as stated, are false for
them. Both failures follow from results in the draft.

\begin{enumerate}
\item \emph{(LH1) fails for the $\pm$odd classes.} The sum
$\Theta_{\pi_0^+}+\Theta_{\pi_0^-}$ is a principal series character and
vanishes on the elliptic set. Hence
$\tr\pi_0^+(f)=-\tr\pi_0^-(f)$ for every $f$ with vanishing hyperbolic
orbital integrals. On the $+$odd class neither vanishes identically
(Lemma~5.29(b) with $k=3$), so two inequivalent unitary representations
give linearly dependent functionals.
\item \emph{(LH1) fails for the even class.} On the elliptic set
$\Theta_{\pi_{nt}}+\Theta_{\pi_1}+\Theta_{\pi_{-1}}=0$ (third line of
(5.18)). More generally
$\Theta_{\pi_m}+\Theta_{\pi_{-m}}+\Theta_{F_m}=0$ there, where $F_m$ is the
$m$-dimensional representation. (LH1) quantifies over all irreducible
admissible representations, so these relations violate it as soon as $L$
contains enough $K$-types.
\item \emph{(LH2) fails for the even class.} By the lemma on non-tempered
representations (Lemmas 5.44 and 5.45),
$\tr\ov\pi_{+,nt}(\ov f)=-\tr\pi_{-1}(f)$ on the even class. On the
elliptic set, the irreducible characters of $SL(2,\R)$ times
$e^{i\theta}-e^{-i\theta}$ are $-\sgn(m)e^{im\theta}$, $\mp1$,
$e^{im\theta}-e^{-im\theta}$ and $0$. None equals $-e^{-i\theta}$, so
$t_\infty(\ov\pi_{+,nt})$ does not exist, while the trace does not vanish
on the class.
\item \emph{(LH3) and $K$-types.} If $f$ has the single $K$-type $12$ and
vanishing hyperbolic orbital integrals, then
$\tr\pi_{11}(f)=-\tr\pi_{-11}(f)-\tr F_{11}(f)=0$, because neither
$\pi_{-11}$ nor $F_{11}$ contains the type $12$. Detection of
$\pi_{\pm11}$ therefore needs an index $L$ containing $-12$ or a $K$-type
of $F_{11}$. This should be said in the final example of \S7.4.
\end{enumerate}

The same relations hold at $p$-adic places: on functions with vanishing
hyperbolic orbital integrals, the two constituents of a reducible unitary
principal series have opposite traces, and the trivial and Steinberg
characters are negatives of each other. So no class of the type in
Definition~5.18 can satisfy (LH1) as stated.

\begin{fix}[Suggested repair of Hypothesis 7.3]
The proof of Theorem~7.9 uses (LH1) only for the representations in
$\Lambda_v$, and (LH2) only for local components of representations in
${\cal A}_{00}$. Fix a set ${\cal F}_v$ of irreducible unitary
representations of $G(k_v)$ containing every local component of a cuspidal
automorphic representation, and a set $\ov{\cal F}_v$ of genuine ones
containing every local component of a representation in ${\cal A}_{00}$.
Require (LH1) only for $\sigma_j\in{\cal F}_v$, and (LH2) only for
$\ov\tau\in\ov{\cal F}_v$ with $t_v(\ov\tau)\in{\cal F}_v$. (The even Weil
representations are not in $\ov{\cal F}_v$, since their Waldspurger lifts
are one-dimensional.) Definition~7.4 must then define ``lies over $\nu$''
only for $\ov\pi'$ in ${\cal A}_{00}$, and Hypothesis~GH must be stated
for the ${\cal A}_{00}$ part of the spectrum.
\end{fix}

Even with this repair, item~1 shows that the $\pm$odd classes cannot be
used when $\pi_\infty$ is a limit of discrete series. A further small
point: (LH2) does not say that $t_v(\ov\tau)$ is unique. It is, by (LH1),
and one sentence should say so.

\subsection*{M4. Remark 2026.1 is false when \texorpdfstring{$q\equiv3\pmod4$}{q = 3 mod 4}}

Remark~4.2 (\texttt{rem:sc-caveat}) says that the function
$(x,\zeta)\mapsto\zeta\chi_\pi(x)$ ``is not obviously invariant under
conjugation.'' It is not invariant, and this refutes relation (4.4).

Let $q\equiv3\pmod4$, $e\in{\cal O}^\times$, and
\[
h=\begin{pmatrix}-1+\varpi e&-e\\ \varpi&-1\end{pmatrix}\in K_0,\qquad
w=\begin{pmatrix}0&-1\\1&0\end{pmatrix}.
\]
Then $\det h=1$, and $(\tr h)^2-4=-4\varpi e+\varpi^2e^2$ has valuation
$1$, so $h$ is regular elliptic (ramified). A direct computation with
(\ref{eq:conjfactor}) below gives
\[
(w,1)\,(h,1)\,(w,1)^{-1}=(whw^{-1},-1).
\]
In general, for $g=\left(\begin{smallmatrix}a&b\\c&d\end{smallmatrix}\right)$
with $abc\neq0$,
\begin{equation}\label{eq:conjfactor}
(w,1)(g,1)(w,1)^{-1}=\bigl(wgw^{-1},\,(a,-bc)_2\bigr),
\end{equation}
and here $(a,-bc)_2=(-1+\varpi e,\varpi e)_2=(-1,\varpi)_2=-1$.
Equivalently $s_v(h)=-1$ and $s_v(whw^{-1})=1$. If (4.4) held, then
\[
\chi_\pi(h)=\chi_{\ov\pi}(h,1)=\chi_{\ov\pi}(whw^{-1},-1)
=-\chi_\pi(whw^{-1})=-\chi_\pi(h),
\]
so $\chi_\pi(h)=0$ for every supercuspidal $\pi$. That is false. The reduction of $h$ is $-1$ times a regular unipotent,
and the same holds at the adjacent vertex $x_1$, so the fixed point set of
$h$ in the tree is the single edge $\{x_0,x_1\}$. For
$\pi=\text{c-Ind}_{K_0}^{G}\sigma$ of depth zero the Frobenius formula
(Proposition~4.4 for $G$) therefore has one term,
$\chi_\pi(h)=\chi_\sigma(\bar h)$, and the cuspidal characters of
$SL(2,\mathbb{F}_q)$ do not vanish at $-1$ times a regular unipotent. (The
sign computation is verified by machine for $q=3,7,11$; for $q=5,13$ the
sign is $+1$, as expected.)

Two consequences for the text.
\begin{itemize}
\item The parenthesis in Remark~4.2 should say that $s_v$ is nontrivial on
regular elliptic elements of $K_0$ \emph{if and only if}
$q\equiv3\pmod4$. When $v(c)$ is odd and $g$ is elliptic, $d\equiv\pm1$
modulo $\varpi$, so $s_v(g)=(\varpi,d)_2=1$ unless $-1$ is a non-square.
\item ``A corrected statement presumably involves the factor $s_i$'' is
optimistic. In the Frobenius formula the sign
$s_i(gxg^{-1})$ varies with the coset representative $g$, so the genuine
character is a \emph{signed} sum of the terms of $\chi_\pi$, with no
evident relation to $\chi_\pi$ itself. I suggest relabeling
Remarks~2026.1 and 2026.2 as ``false when $q\equiv3\pmod4$; open when
$q\equiv1\pmod4$,'' or removing Chapter~4 from the main text
(\S\ref{sec:read}).
\end{itemize}

\subsection*{M5. The weighted orbital integral at central elements diverges}

Section~5.2.5 \lref{1503--1540} and the example in \S6.1.1
\lref{3615--3646} define $J_A(z,f)$, $z$ central, as
$\lim_{a\to1}J_A(az,f)$, and state $r_A(\gamma,a)=1$. The limit does not
exist in general. For $a=a(t)$ and $x=n(u)k$ one has
$x^{-1}ax=k^{-1}a\,n\bigl((1-t^{-2})u\bigr)k$, so in the $p$-adic case
\[
J_A(a,f)=|t|\int_K\!\int_F f\bigl(k^{-1}a\,n(u)k\bigr)\,
\log\max\!\Bigl(1,\frac{|u|}{|1-t^{-2}|}\Bigr)\,du\,dk
=-\log|1-t^{-2}|\cdot J_G(a,f)+O(1),
\]
which diverges as $t\to1$ unless the unipotent orbital integral of $f$
vanishes. Arthur's definition \cite[(8.1)]{A7} contains the term $L=G$
with a nonzero function $r_A^G(z,a)$ precisely to cancel this.

\begin{fix}[Suggested replacement for the last display of \S5.2.5]
For $z=\pm1$ and $a=a(t)\to1$,
\[
J_A(z,f)=\lim_{t\to1}\Bigl[J_A(az,f)+\log|1-t^{-2}|\,J_G(az,f)\Bigr]
=\int_K\!\int_F f\bigl(k^{-1}z\,n(u)k\bigr)\,\log|u|\,du\,dk
\]
in the $p$-adic case, with $\log|u|$ replaced by the corresponding
archimedean weight otherwise. (The constant in front of the logarithm
depends on the normalization of $v_A$; check it against \cite{A7}.)
\end{fix}

The same display has $dx/dt$ where $dx/da$ is meant, and the displays of
\S5.2.1 mix the variables $g$, $a$ and $\gamma$.

\subsection*{M6. Chapter 6 does not state the trace formula correctly}

The chapter is not used in the proof of Theorem~7.9, but Remark~7.8
(\texttt{rem:GH}) builds on it, and it is presented as a recollection of
known results. As written:
\begin{enumerate}
\item The spectral side omits the discrete terms that are not cuspidal:
the residual spectrum (the trivial representation of $SL(2)$; the even
theta representations of the cover) and the terms
$\tfrac14\tr\bigl(M(w,\chi)I_B(\chi,f)\bigr)$ for $\chi^w=\chi$. The lemma
\lref{3732--3757} treats only ``unramified'' $\chi$, that is
$\mu^w\neq\mu$. In a comparison these terms cannot be ignored: the lemma on
non-tempered representations in \S5.7.4 is the local shadow of the
matching of residual spectra.
\item In (e), the integral should be over automorphic characters of
$A(k)\backslash A(\A)^1$ (times $i\mathfrak{a}^*$), not over
$\Pi_u(A(\A_S))$. The operator is written $M_A^T$ in (e) but $M_B^T$ in
the definition.
\item The formula depends on the truncation parameter $T$ through
$-2T\cdot I$. The identity is stated for $SL(2)$ without fixing $T$ and for
the cover at $T=0$.
\item In the definition \lref{3705--3730}, the clause
``if $\chi=(G,\rho)$, let $M_B^T(\pi)_\chi=1$'' duplicates the previous
one and makes no sense for $P=B$. The factor
$\operatorname{vol}(\R/\Z)$ is unexplained.
\item On the geometric side, the sums run over
$\gamma\in{\cal O}\cap A(\A)$ where $A(k)$ is meant, and
$a^M(S,\gamma)$ is given as
$\operatorname{vol}(M(k)_\gamma\backslash M(\A_S)^1_\gamma)$ where the
global volume $\operatorname{vol}(M(k)_\gamma\backslash M(\A)^1_\gamma)$ is
meant. The display \lref{3589} has unbalanced delimiters.
\item The two displayed trace formulas in Remark~7.8(1) inherit the
omission in item~1, and reuse the letter ${\cal C}$ for the set of classes
${\cal O}$.
\end{enumerate}
I suggest replacing Chapter~6 by a one-page statement of the two formulas
with precise references, or moving it to an appendix marked ``1995 text,
not used.''

\subsection*{M7. Three incorrect statements in proved or explanatory text}

\begin{enumerate}
\item \emph{Remark 2026.11} (\texttt{rem26:A00-norm}) says that any bound
strictly better than $|s_v|<1/2$ would do. It would not. The proof of
Lemma~7.12 needs $2s_v\neq\pm1/2$, and a bound such as $|s_v|\le 1/3$ does
not exclude $s_v=1/4$. What is needed is $|s_v|\neq1/4$ at one suitable
place, and the Gelbart--Jacquet bound $|s_v|<1/4$ is exactly the natural
sufficient condition.
\begin{fix}
The proof of Lemma~7.12 uses the bound $|s_v|<1/4$ only to exclude
$|s_v|=1/4$ at a single place $v\notin S$ where all data are unramified.
The trivial bound $|s_v|<1/2$ does not suffice.
\end{fix}
(Under the even correspondence of M1 the theta representations match
$s=1$ on $SL(2)$, and the trivial bound would suffice.)

\item \emph{Lemma 3.4} \lref{951--958}. Since $J_G$ already
contains the factor $|D(\gamma)|^{1/2}=|a_1-a_2|$, the conclusion should
read $F^\vee(\gamma)=J_G(\gamma,f)$. The extra factor $|a_1-a_2|$ counts
the normalization twice. The letter $F$ also clashes with the field.

\item \emph{Lemma 3.2} \lref{893--915}. The hypothesis
$\chi\neq\chi^w$ does not imply that $I_A(\chi)$ is irreducible (take
$|x|^{\pm1}$ on $SL(2)$ or $|x|^{\pm1/2}$ on the cover), and it is not
needed: the formula is the character of the induced representation for
every $\chi$. Proposition~5.12 applies the formula to all $\chi$, while
\S7.4 item~2 restricts to $\chi\neq\chi^w$. State the lemma for all
$\chi$, as the character of the (possibly reducible) representation
$I_A(\chi)$, and delete the restriction in \S7.4.
\end{enumerate}

\section{What is correct}

So that the list above is read in proportion, here is what I verified.

\begin{itemize}
\item \S2.1: the two forms of the cocycle (2.1) agree; (2.3)--(2.5);
Lemma~2.2; Lemma~2.3(a), (e), (f), (h). For (h) I found a clean proof,
which could replace the appeal to the literature: for $c\neq0$,
$\beta(g,n)=(1,c)=1$, $\beta(g,g^{-1})=(c,-c)=1$, and
$\beta(gn,g^{-1})=(-c,c)=1$ because $gng^{-1}$ has lower left entry
$-c^2$.
\item \S2.2: the quotient construction, $Z_0$, and the splitting over
$G(k)$ are consistent. (In the $\beta_{\A}$ model the image of
$\gamma\in G(k)$ is $(\gamma,s_{\A}(\gamma))$; the inverse in the text is
harmless since $s_{\A}=\pm1$.)
\item \S5.3.1--5.3.2: Lemmas 5.3, 5.4, 5.8, 5.15, Definitions 5.5 and 5.7, and both
$\Q_3$ examples. The correction of the 1995 requirement
$\epsilon=\vartheta$ is right.
\item \S5.3.4: the transfer of $K\cap A$-invariance under Definition~5.7; the
unit elements (Example~5.22).
\item \S5.5: formulas (5.17); the $SL(2,\R)$ character formulas and central
characters; the Weyl integration formulas with $d\theta/2\pi$ and
$d\varphi/4\pi$; Remark~5.27; the example following it; the table and all
three parts of Lemma~5.29; the infinitesimal character bookkeeping
$|m|\in\{2\nu-1,2\nu,2\nu+1\}$.
\item \S5.4: items 1--6 of the list of errors in the 1995 dyadic argument;
the $\Q_2$ example.
\item Lemma~5.36 (Fresnel integral).
\item Chapter 7: Lemma~7.10 (given Harish-Chandra's theorem, and with the
draft's own caveat for the cover); Lemma~7.11; Lemma~7.12, including the
Jacquet module and the normalization $\lambda=\omega|\cdot|^{-1/2}$;
Steps~1--4 of the proof of Theorem~7.9; Proposition~7.15;
\S7.5, both items.
\end{itemize}

\section{Minor mathematical points}

\begin{enumerate}
\item \emph{Lemma 2.3(b)} \lref{544--546}. For $p$ odd the splitting of
$\ov{K_0}$ is unique, because $SL(2,{\cal O})$ has no quotient of order
$2$. So ``a splitting'' can be ``the splitting,'' and
$\kappa_v(k)=(k,s_v(k))$ holds at \emph{every} odd place, not only almost
all \lref{700}. Uniqueness also settles a point left implicit in
Proposition~7.15: conjugation by $w$ in the cover of $GL(2)$ carries
$\kappa(K_0)$ to the splitting of $\ov{K_1}$.
\item \emph{The cocycle on $GL(2)$} \lref{473--482}. The definition sets
$v=1$ when $F=\R$. My recollection of \cite[\S2.1]{Ge2} is that
$v(y,s)=(y,d)_2$ for $c=0$ with no exception at the real place. Please
check.
\item \emph{\S3.1.2} \lref{886--891}. For regular $t$ in $SL(2)$ the
centralizer \emph{is} $T$; ``finite index'' is not wrong but misleading.
\item \emph{\S4.2} \lref{1111--1116}. The quotient
``$(N\cap K_{i,m-1-i})/K_{i,m}$'' is not a quotient of groups as written.
Presumably $(N\cap K_{i,m-1})K_{i,m}/K_{i,m}$ is meant; the index
$m-1-i$ needs an explanation.
\item \emph{Lemma 4.5} \lref{1225--1228}. The
quantifiers are wrong. It should read: there are a neighborhood $U$ of $x$
and a finite set $Z$ such that $\chi^z_{\ov\sigma}(y)=0$ for all $y\in U$
and all $z\notin Z$. In the proof, ``unless $p(x)$ is elliptic regular''
should be ``unless $p(x)$ is elliptic'' ($x$ is regular by assumption).
\item \emph{Remark 5.30} \lref{2675--2681}. The factor
$\sgn(x)^{-\varepsilon k/2}$ is ambiguous for $x<0$ and half-integral
exponent, and $\ln(x)$ should be $\ln|x|$. Write the fourth root of unity
explicitly. The final sentence of the remark asserts an inconsistency
without showing it.
\item \emph{(5.18)} \lref{3383--3389}. Lines 3 and 4 give the same
decomposition for two different principal series,
$\pi_{PS}(0,1)$ and $\pi_{PS}(0,-1)$. Presumably the parameters are the
exponents $\pm1$. Adams's notation should be explained.
\item \emph{\S5.7 is used later.} The section says ``This section shall not
be needed later'' \lref{2817}, but formula (iii) of \S5.7.2 is used in
Lemma~7.12, and \S5.7.4 in \S7.4 and M3.
\item \emph{\S5.7.2--5.7.3}, all conditional on the Conjecture:
\begin{itemize}
\item \lref{3049--3056}: the kernel $k_\epsilon(x,y)$ is written with
$\tfrac{a+d}{c}q(x)$ and then with $\tfrac{a+d\mp2}{c}xy$. Off the
diagonal the phase is $(ax^2+dy^2\mp2xy)q/c$; the displayed formulas are
right only for $x=y$. A factor $\tfrac12$ is lost between the kernel and
the following Remark.
\item \lref{3254--3256}: $|\det(g-1)|^{1/2}=2|\sin(\theta/2)|$, so the
second line should have $2|\sin(\theta/2)|$ and $2|\cos(\theta/2)|$,
without the exponent $1/2$. (This is what agrees with the formulas quoted
from Adams.)
\item $\gamma(\tau)=\exp(\tfrac{\pi i}{4}\sgn q)$ \lref{3166} and
$\gamma(q)=\exp(-\tfrac{i\pi}{4}\sgn q)$ \lref{3265} contradict each
other.
\item \lref{3123}: the $n$-dimensional Fresnel integral is missing the
phase $e^{\pi i n\,\sgn(a)/4}$, and $F^n$ should be $\R^n$.
\item \lref{3029--3031}: the ``trick'' ignores the scalar $\omega(-1)$ by
which $(-I,1)$ acts and the cocycle $\beta(-I,g)$.
\item \lref{3007}: $L^2(G)_\epsilon$ should be $L^2(F)_\epsilon$;
\lref{3232}: stray ``$|$''.
\end{itemize}
\item \emph{Lemma 5.23} \lref{2206--2213} begins
``Finally'' and repeats Example~5.22. Merge them.
\item \emph{Hypothesis 7.3 (LH0) and Step 1.} ``The corresponding genuine
$\ov{K_v}$-types'' in (LH0) is not defined. Since $\ov K=\ov{SO(2)}$ has
characters $e^{in\theta/2}$, say which half-integers correspond to a given
finite set $L$.
\end{enumerate}

\section{Readability}\label{sec:read}

\begin{enumerate}
\item \emph{The main line is hard to find.} Of roughly 70 pages, the main
theorem uses \S2, \S5.3, formula (iii) of \S5.7.2, and Chapter~7. Chapters
3, 4 and 6, and \S\S5.4, 5.6, 5.7.3--5.7.4, are either unverified 1995
material or unused. I suggest three parts: (I) background and the verified
local results; (II) the conditional theorem; (III) an appendix,
clearly marked, holding the unverified and unused 1995 material. A
dependency diagram in the introduction would help.

\item \emph{Discussion of the 1995 errors dominates the text.} The
introduction, \S5.4, \S7.5 and many remarks are about what the old draft
got wrong. For a reader who never saw the 1995 draft this is a
distraction. In line with the convention the author has adopted for other
revisions, this material belongs in the revision log, with any
mathematically useful residue rephrased as a remark (for instance,
Remark~5.11 on why weighted orbital integrals cannot match for all pairs
stands on its own without mentioning 1995).

\item \emph{The labeling convention of \S1.3 is not followed.} It says
that statements labeled Theorem, Proposition, Lemma or Corollary are proved
or quoted with a reference. Exceptions: the lemma on the existence of
characters on the cover (\S3.1.1, ``no proof has been published''; cite
Li's local harmonic analysis paper, \S\ref{sec:bib}); the principal series
character lemma (``proof omitted''); Lemma~7.10 for the cover (flagged in
Remark~2026.10); Proposition~5.38 and Corollaries 5.39 and 5.42, which
assume a Conjecture; Lemma~6.3.

\item \emph{``Remark 2026.$n$'' is a confusing name.} These are not
remarks; they are unverified or false \emph{claims}. Meanwhile the
critiques of those claims are ordinary Remarks (4.2, 5.11, 5.31). Rename the
environment, for example ``Unverified claim 2026.$n$,'' so that a reader
skimming for statements is not misled.

\item \emph{Notation that is used before it is defined, or never
defined:} $\Pi(\cdot)$, $\Pi_u(\cdot)$, $\Pi_{gen}$; $I_A(\chi)$ (the
normalization of induction is never stated); $J_G$, $F_f^{\ov{A(F)}}$ and
$F^\vee$ in \S3.1.2 (defined in \S5.2 or not at all); $\overline C$ and
``the argument of \cite{Ar} sketched above'' in \S3.2.2 (nothing is
sketched above); $t$ and $\Delta$ ``as in \cite{Fl}''; $\omega$, $q(x)$,
$\gamma(q)$ ``as in \cite{Fl}.''

\item \emph{Overloaded symbols.} The following each carry several
meanings, sometimes on the same page:

\medskip
\begin{center}
\begin{tabular}{@{}ll@{}}
\toprule
Symbol & Meanings\\
\midrule
$w$ & $\left(\begin{smallmatrix}0&-1\\1&0\end{smallmatrix}\right)$;
$\left(\begin{smallmatrix}0&\varpi\\1&0\end{smallmatrix}\right)$ (\S4.2, Ch.~7)\\
$p$ & residual characteristic; projection $\ov G\to G$; map $GL_2\to SL_2$ (also $pr$)\\
$\nu$ & valuation; infinitesimal character; $k/2-1$\\
$\epsilon,\varepsilon$ & genuine parameter; sign $\pm1$ of $\chi_\epsilon$ on $\R$; parity; identity character of $\mu_2$\\
$\tau$ & map $h\mapsto(h,1)^2$; $\pm1/2$; additive character; Ramanujan $\tau(p)$\\
$\psi$ & additive character; the function in Lemma~5.4; elements of $\Psi(P_1,P,\chi,W)$\\
$\chi$ & character of $A$; class in ${\cal X}$; quadratic character\\
$S$ & set of places; Satake transform\\
$B$ & Borel subgroup; Iwahori subgroup (\S5.6.1)\\
$T$ & Cartan subgroup; truncation parameter\\
$q$ & residue field order; quadratic form $q(x)$\\
$F$ & local field; the function $\chi\mapsto\tr I_A(\chi)(f)$\\
${\cal A}$, ${\cal C}$ & finite sets in Step~1 and spaces ${\cal A}_0$; classes ${\cal C}_v$ and set of ${\cal O}$\\
\bottomrule
\end{tabular}
\end{center}
\medskip

The notation table lists both $a(x)$ and $h(x)$ for the same matrix; pick
one.

\item \emph{Examples placed at chapter ends.} Several examples sit after
the last section of a chapter and refer to ``the lemma above'' many pages
back (end of Chapters 3, 4, 5, 6). Move each next to the statement it
illustrates.

\item \emph{Small things.} ``In the case case'' \lref{3736}; ``a
irreducible'' \lref{1125}; ``non-metaplectic Case'' \lref{2119}; ``Define
$I_P(\pi)$ induced representation \dots\ Define \dots\ denote''
\lref{3699--3703}; $\operatorname{Ind}$ for c-Ind in Remark~2026.1;
``$cos$'' unformatted \lref{2853}; ``$ln$'' and ``$<<$'' in the proof of
Lemma~5.36; ``p-adic'' versus ``$p$-adic''; first person singular and
plural are mixed (``Let me first prove,'' ``we obtain'').
\end{enumerate}

\section{Literature and bibliography}\label{sec:bib}

\begin{enumerate}
\item \emph{Chai--Zhang} \lref{99--102, 260--265}. Their theorem also
requires the two representations to be generic with respect to the same
additive character. Add this condition in both places.
\item \emph{[Zh]} is Q.~Zhang, ``Addendum to `A strong multiplicity one
theorem for $SL_2$','' Pacific J. Math. 292 (2018), no.~1, starting at
p.~283. The bibliography gives the title of the original paper and
``283ff.''
\item \emph{Add:} W.-W.~Li, ``Stabilization of the trace formula for
metaplectic groups,'' arXiv:2109.06581; W.-W.~Li, ``Spectral transfer for
metaplectic groups. I. Local character relations,'' J. Inst. Math. Jussieu
18 (2019) 25--123. From memory, to be verified: W.-W.~Li, ``La formule des
traces pour les rev\^etements de groupes r\'eductifs connexes. II. Analyse
harmonique locale,'' Ann. Sci. \'Ec. Norm. Sup\'er. 45 (2012) (existence of
characters on covering groups, for \S3.1.1); W.~T.~Gan and G.~Savin,
``Representations of metaplectic groups II: Hecke algebra
correspondences,'' Represent. Theory 16 (2012) (cited by name in \S5.4 but
absent from the bibliography); H.~Jacquet and J.~Shalika, ``On Euler
products and the classification of automorphic representations I, II,''
Amer. J. Math. 103 (1981) (for M1).
\item \emph{[LL]}: Canad. J. Math. \textbf{31} (1979) 726--785, not volume
33 (from memory).
\item \emph{[M1], [M2]}: Math. Ann. 266 appeared in 1984, not 1987 (from
memory).
\item \emph{[J1]}: I believe this is Israel J. Math. \textbf{61} (1988)
113--154; the entry says volume 66. The author will know.
\item \emph{[Ad1]} has appeared: Israel J. Math. 98 (1997) 229--252 (from
memory).
\item \emph{Author dashes.} The entry [A1] begins ``------'' directly
after Ariturk, so it reads as a paper by Ariturk; write ``J.~Arthur.'' The
same problem occurs at [GS] (after Gelbart--Piatetski-Shapiro) and should
be checked throughout.
\item \emph{Uncited entries:} [AC], [Bl], [Bo], [HC1], [L2], [Sa], [Sch].
Remove or cite.
\end{enumerate}

\section{Options for the revision}\label{sec:options}

\begin{enumerate}
\item[(A)] \emph{Refocus on the local results.} Present the paper as a
study of the genuine characters of the torus, the matching of functions on
it, and the real elliptic matchings, with Lemma~7.12 as an application.
State Proposition~R as a theorem: it is a clean negative result explaining
why the odd unramified correspondence has no global counterpart. This is
the smallest change and yields a correct paper.

\item[(B)] \emph{Rebuild the conditional theorem on the even
correspondence.} Replace Definition~5.21 by the homomorphism
$b:{\cal H}(G,K)\to{\cal H}(\ov G,K)_{gen}$ with
$\tr I(\chi_\psi\mu)(b(f))=\tr I_A(\mu^2)(f)$. This is canonical and
agrees with the archimedean matchings already in the draft. But the
argument changes in essential ways, and I do not see how it yields
multiplicity one:
\begin{itemize}
\item $b$ is not an isomorphism and the transfer is not injective: $\mu$
and $\mu$ times the unramified quadratic character have the same image. In
Step~2 of the proof the right side becomes a sum over all $\ov\pi'$ whose
Waldspurger lifts are quadratic twists of one another, and Step~3
(uniqueness from injectivity of $t_v$) is lost.
\item Only those $\pi$ are reached whose $\Pi$ can be twisted to have
trivial central character. For the others no $\ov\pi'$ exists, so a trace
identity of the form GH cannot hold on the whole cuspidal spectrum.
\item The fibers are Waldspurger packets on one side and
Labesse--Langlands packets on the other. An identity relating them will
be of endoscopic type, with signs, not an equality of cuspidal traces.
\end{itemize}
If the author pursues this, the honest statement of the result would be a
description of how the two packet structures correspond, with multiplicity
one for $SL(2)$ as an input (Ramakrishnan), not an output.

\item[(C)] \emph{Keep Theorem 7.9 but say what it is.} If the theorem is
retained for the record, the abstract and introduction must state that its
hypotheses are inconsistent (Proposition~R), so that it carries no
information about multiplicity one. I do not recommend this.
\end{enumerate}

In all three cases the abstract's sentence ``we show that $\pi$ occurs with
multiplicity one and corresponds to a unique genuine cuspidal automorphic
representation of the cover'' must go.

\begin{thebibliography}{9}
\bibitem[A4]{A7} J. Arthur, ``The local behaviour of weighted orbital
integrals,'' Duke Math. J. 56 (1988) 223--293.
\bibitem[Ge2]{Ge2} S. Gelbart, \emph{Weil's representation and the
spectrum of the metaplectic group}, Lecture Notes in Math. 530, Springer,
1976.
\bibitem[GJ2]{GJ78} S. Gelbart and H. Jacquet, ``A relation between
automorphic representations of $GL(2)$ and $GL(3)$,'' Ann. Sci. \'Ec. Norm.
Sup. 11 (1978) 471--542.
\bibitem[Ar]{Ar} Ariturk, as in the draft.
\bibitem[Fl]{Fl} Y. Flicker, ``Automorphic forms on covering groups of
$GL(2)$,'' Invent. Math. 57 (1980) 119--182.
\bibitem[Wa1]{Wa1} J.-L. Waldspurger, ``Correspondance de Shimura,''
J. Math. Pures Appl. 59 (1980) 1--133.
\bibitem[Wa2]{Wa2} J.-L. Waldspurger, ``Correspondance de Shimura et
quaternions,'' Forum Math. 3 (1991) 219--307.
\end{thebibliography}

\end{document}
